\chapter{Thyroid Antibodies}

\section*{What Is This Test?}
Thyroid antibody tests measure autoantibodies directed against thyroid-specific antigens that are markers of autoimmune thyroid disease. The three principal thyroid autoantibodies tested are: (1) Anti-thyroid peroxidase antibodies (anti-TPO), directed against thyroid peroxidase, the enzyme responsible for iodide oxidation and thyroglobulin iodination in thyroid hormone synthesis --- these are the most clinically useful thyroid antibodies and are present in 90--95\% of Hashimoto thyroiditis and 50--80\% of Graves disease; (2) Anti-thyroglobulin antibodies (anti-Tg), directed against thyroglobulin, the 660 kDa glycoprotein precursor of T3 and T4 stored in thyroid follicular colloid --- present in 60--80\% of Hashimoto thyroiditis and 25--50\% of Graves disease, and important as an interferent in thyroglobulin tumor marker assays; (3) TSH receptor antibodies (TRAb), divided into thyroid-stimulating immunoglobulins (TSI) that activate the TSH receptor causing Graves hyperthyroidism, and TSH receptor-blocking antibodies (TBAb) that block TSH action causing hypothyroidism in a subset of patients with atrophic thyroiditis. Thyroid antibody testing confirms the autoimmune etiology of thyroid dysfunction, predicts progression from subclinical to overt disease, and guides treatment decisions.

\section*{Alternative Names}
Anti-Thyroid Peroxidase Antibodies (Anti-TPO, TPOAb); Anti-Thyroglobulin Antibodies (Anti-Tg, TgAb); TSH Receptor Antibodies (TRAb); Thyroid-Stimulating Immunoglobulin (TSI); TSH-Binding Inhibitory Immunoglobulin (TBII); Anti-Microsomal Antibodies (historical term for anti-TPO); Thyroid Autoantibody Panel

\section*{Principle}
Anti-TPO and anti-Tg antibodies are measured using competitive or non-competitive immunoassays. Most commonly, automated chemiluminescent immunoassays are used: Roche Cobas ECLIA (electrochemiluminescence), Abbott Architect CMIA, Siemens IMMULITE/Atellica, or Beckman Coulter DxI. In competitive immunoassays, patient antibodies compete with labeled anti-TPO antibodies for binding to TPO antigen-coated solid phase. Results are reported in IU/mL calibrated against the WHO International Reference Preparation MRC 66/387 (anti-TPO) and MRC 65/93 (anti-Tg). TSH receptor antibodies (TRAb) are measured using two main methods: (1) competitive binding assay (TBII --- TSH-binding inhibitory immunoglobulin) which measures the inhibition of labeled TSH binding to recombinant TSH receptors by patient immunoglobulins, expressed as \% inhibition or IU/L, calibrated against WHO standard NIBSC 90/672; (2) bioassay measuring cyclic AMP generation in cell lines expressing TSH receptor after stimulation by patient immunoglobulins --- this assay differentiates stimulating (TSI) from blocking (TBAb) antibodies, reported as \% of baseline. Current automated TRAb platforms (Roche Cobas ECLIA, Siemens IMMULITE) use a liquid-phase competitive binding format with recombinant human TSH receptor and highly purified bovine TSH.

\section*{History}
The concept of thyroid autoimmunity was first proposed in 1956 when Deborah Doniach and Ivan Roitt discovered autoantibodies in Hashimoto disease patients, and by Duncan Adams and H.D. Purves who discovered the long-acting thyroid stimulator (LATS) in Graves disease sera. The "thyroid microsomal antigen" was described in the 1960s and eventually identified as thyroid peroxidase (TPO) in 1985 by Czarnocka and colleagues. The WHO International Reference Preparations for anti-TPO (MRC 66/387) and anti-Tg (MRC 65/93) were established in 1993. TSH receptor antibodies were first measured by bioassay in the 1960s (mouse bioassay for LATS). Competitive TBII assays were developed in the 1980s, and automated TRAb assays were introduced in the early 2000s. The 2017 ATA pregnancy guidelines recommend TRAb measurement in pregnant women with Graves disease. The understanding that anti-TPO positivity predicts progression from subclinical to overt hypothyroidism (2--5\% per year) was established by the Whickham Survey follow-up study (Vanderpump et al., 1995).

\section*{Why This Test Is Ordered}
\begin{itemize}
  \item Confirming autoimmune etiology in primary hypothyroidism --- elevated anti-TPO in 90--95\% of Hashimoto thyroiditis, anti-Tg in 60--80\%
  \item Supporting the diagnosis of Graves disease --- anti-TPO positive in 50--80\%, TRAb positive in >95\%
  \item Predicting progression from subclinical hypothyroidism (TSH 4.5--10 mIU/L, normal FT4) to overt hypothyroidism --- anti-TPO-positive patients progress at 4.3\% per year vs. 1.3\% in antibody-negative patients
  \item Predicting postpartum thyroiditis in anti-TPO-positive pregnant women (33--50\% develop postpartum thyroid dysfunction)
  \item Predicting risk of reproductive failure --- anti-TPO positivity is associated with increased miscarriage risk (2--4 fold) and preterm delivery
  \item Detecting anti-Tg antibodies as an interferent in thyroglobulin tumor marker assays --- anti-Tg can cause falsely low or undetectable thyroglobulin in differentiated thyroid cancer follow-up
  \item Predicting amiodarone-induced thyroid dysfunction --- anti-TPO-positive patients have higher risk when starting amiodarone
  \item Evaluating patients receiving interferon-alpha, interleukin-2, or immune checkpoint inhibitors (pembrolizumab, ipilimumab) who develop thyroid dysfunction --- drug-induced destructive thyroiditis vs. autoimmune activation
  \item Predicting risk of fetal/neonatal hyperthyroidism in pregnant women with active or past Graves disease --- TRAb >3 times upper limit of normal in second/third trimester indicates significant risk of neonatal Graves disease
  \item Investigating goiter etiology --- anti-TPO frequently positive in goitrous Hashimoto thyroiditis
  \item Diagnosis of IgG4-related thyroid disease (Riedel thyroiditis, fibrosing variant of Hashimoto) --- associated with elevated IgG4 and anti-TPO
\end{itemize}

\section*{Primary vs Secondary Test}
This is a secondary (adjunctive) test. Thyroid antibody testing is used to confirm the autoimmune etiology of abnormal thyroid function detected by TSH screening, to predict the course of disease, and to guide management decisions. Anti-TPO is recommended by ATA guidelines when TSH is elevated (particularly 4.5--10 mIU/L) to predict progression risk. TRAb is recommended in all pregnant women with Graves disease (active or past) by the 2017 ATA pregnancy guidelines. Routine screening of asymptomatic individuals with thyroid antibodies is not recommended.

\section*{How This Test Is Performed}
A healthcare professional draws a blood sample from a vein in your arm into a serum separator tube (gold-top) or plain red-top tube. Approximately 5 mL of blood is collected. The specimen is centrifuged to separate serum. Anti-TPO, anti-Tg, and TRAb are stable in serum at 2--8 degrees C for 7 days and at -20 degrees C for 6 months. Frozen samples should be thawed completely and mixed before testing. Avoid repeated freeze-thaw cycles. Separate samples are typically used for each antibody test, though some laboratories offer a thyroid antibody panel.

\section*{What Preparation Is Needed}
No special preparation is required. Fasting is not necessary. Thyroid antibodies do not show significant circadian variation. High-dose biotin (>5 mg/day) should be discontinued 48 hours before testing because biotin interferes with streptavidin-biotin immunoassay platforms, potentially causing falsely elevated TRAb or anti-TPO results (biotin binds streptavidin, reducing signal generation). For TRAb measurement in pregnancy, testing is recommended at 20--24 weeks gestation in women with active or prior Graves disease. Patients receiving intravenous immunoglobulin (IVIG) may have transiently positive thyroid antibodies from donor immunoglobulins. Therapeutic monoclonal antibodies (e.g., rituximab) can interfere with some immunoassays.

\section*{Contraindications and Precautions}
\begin{itemize}
  \item No absolute contraindications. Important considerations: anti-TPO antibodies can be detected in 10--15\% of euthyroid healthy individuals (more common in women and with advancing age), so a positive result without TSH elevation does not necessarily indicate current thyroid disease but identifies individuals at future risk. Anti-Tg antibodies are present in 10--15\% of the general population and are less specific than anti-TPO. Anti-Tg antibodies interfere with thyroglobulin measurement in thyroid cancer follow-up --- all thyroglobulin assays must include anti-Tg measurement, and the presence of anti-Tg renders thyroglobulin unreliable as a tumor marker (underestimation on immunometric assays, overestimation on RIA). TRAb levels can be elevated in apparently euthyroid patients with nodular goiter who are at risk of developing hyperthyroidism after iodine load (e.g., amiodarone, contrast agents). Thyroid antibody results cannot distinguish Hashimoto thyroiditis from Graves disease --- clinical presentation, TSH/T4/T3 status, and TRAb must be used for this distinction. Heterophile antibodies and rheumatoid factor can produce false-positive results in some immunoassay platforms.
\end{itemize}
\section*{Risks}
Minimal risk. Slight pain, bruising, or hematoma at the venipuncture site. Rarely: vasovagal syncope, infection, or phlebitis. No radiation exposure.

\section*{What Happens After the Test}
Results are available within 1--4 hours on automated platforms (anti-TPO, anti-Tg) or within 1--3 days for TRAb (if sent to reference laboratory). In a patient with elevated TSH and positive anti-TPO, a diagnosis of Hashimoto thyroiditis is confirmed. If TSH is normal, the patient has "euthyroid Hashimoto thyroiditis" or "antibody-positive euthyroidism" and should have TSH monitored annually (2--5\% per year progression to overt hypothyroidism). Pregnant women with positive anti-TPO should have TSH monitored every 4 weeks during pregnancy and at 3 and 6 months postpartum for postpartum thyroiditis. Patients with Graves disease and positive TRAb require ongoing monitoring and should not receive radioactive iodine therapy if they plan pregnancy within 6 months.

\section*{Understanding Results}

\subsection*{Below Normal}
\begin{itemize}
  \item \textbf{Undetectable/low anti-TPO (<35 IU/mL, lab-specific):} No evidence of autoimmune thyroid disease. However, approximately 5--10\% of Hashimoto thyroiditis is anti-TPO-negative (seronegative autoimmune thyroiditis), confirmed by ultrasound showing hypoechoic heterogeneous parenchyma.
  \item \textbf{Negative TRAb (<1.75 IU/L, lab-specific):} Makes Graves disease unlikely as a cause of hyperthyroidism. Other causes of suppressed TSH should be investigated (toxic nodular goiter, subacute thyroiditis, factitious thyrotoxicosis).
  \item \textbf{Declining TRAb in Graves disease:} Indicates favorable response to anti-thyroid drugs and is one criterion for considering a trial of methimazole cessation after 12--18 months of therapy (relapse rate approximately 30--50\%).
  \item \textbf{Negative thyroid antibodies in subclinical hypothyroidism:} Lower risk of progression (1.3\% per year) compared to antibody-positive (4.3\% per year), and levothyroxine treatment benefit is less certain.
\end{itemize}

\subsection*{Above Normal}
\begin{itemize}
  \item \textbf{Elevated anti-TPO (>35 IU/mL, often markedly elevated >500 IU/mL):} Diagnostic of autoimmune thyroid disease. Most commonly Hashimoto thyroiditis (lymphocytic infiltration of thyroid, progressive follicular destruction, progression to hypothyroidism). Also elevated in 50--80\% of Graves disease. Higher titers correlate with more active lymphocytic infiltration on histology and faster progression to hypothyroidism.
  \item \textbf{Elevated anti-Tg (>20 IU/mL, lab-specific):} Less specific than anti-TPO. Present in 60--80\% of Hashimoto thyroiditis, 25--50\% of Graves disease, and some patients with thyroid cancer. Elevation of either anti-TPO or anti-Tg confirms autoimmune etiology. Anti-Tg is essential to measure alongside thyroglobulin in differentiated thyroid cancer monitoring because anti-Tg interference invalidates thyroglobulin results.
  \item \textbf{Positive TRAb (>1.75 IU/L or >140\% baseline in bioassay):} Diagnostic of Graves disease in a patient with suppressed TSH and hyperthyroidism. TRAb is positive in >95\% of active Graves disease. TRAb >3 times the upper limit of normal in second/third trimester pregnancy is associated with 1--5\% risk of neonatal Graves disease from transplacental transfer. Persistently positive TRAb after 12--18 months of anti-thyroid drug treatment is a predictor of relapse (sensitivity 65\%, specificity 80\%).
  \item \textbf{TSH receptor-blocking antibodies (TBAb):} Present in 5--10\% of Hashimoto thyroiditis, causing atrophic thyroiditis variant (small thyroid with no goiter). Functional bioassay needed to distinguish TSI from TBAb.
  \item \textbf{Anti-TPO positivity with normal TSH (euthyroid):} Identifies autoimmune thyroid diathesis. Risk of progression to overt hypothyroidism is 2--5\% per year. Annual TSH monitoring recommended. Associated with 2--4 fold increased miscarriage risk in euthyroid pregnant women.
  \item \textbf{Amiodarone-induced thyrotoxicosis Type 1:} Positive TRAb and/or elevated anti-TPO in nodular goiter or latent Graves with iodine load-driven hyperthyroidism (Type 1), distinguished from destructive thyroiditis (Type 2) where antibodies are typically absent.
\end{itemize}

\section*{Normal Results}
\begin{itemize}
  \item \textbf{Anti-TPO antibodies:} <35 IU/mL (may vary to <9--35 IU/mL depending on assay; Roche and Abbott negative cutoff commonly <9 IU/mL, others <35 IU/mL)
  \item \textbf{Anti-thyroglobulin antibodies:} <20 IU/mL (lab-specific; Roche <115 IU/mL, Abbott <4.1 IU/mL --- significant inter-assay variability)
  \item \textbf{TRAb (competitive binding assay, TBII):} <1.75 IU/L (Siemens), <1.22 IU/L (Roche) --- results vary significantly by assay
  \item \textbf{TSI bioassay:} <140\% of basal activity (reported by some reference laboratories)
  \item \textbf{Anti-TPO in healthy population:} 10--15\% have elevated anti-TPO without thyroid dysfunction (female preponderance 2:1, prevalence increases with age)
  \item \textbf{Note:} Thyroid antibody assays are not yet harmonized across manufacturers. Use the same assay for serial monitoring because switching platforms can produce step-changes in antibody titer.
\end{itemize}

\section*{Follow-up Tests}
\begin{itemize}
  \item \textbf{TSH and free T4:} Essential to co-interpret with antibody results to determine current thyroid function status and need for levothyroxine therapy. TSH >10 mIU/L with positive anti-TPO indicates overt hypothyroidism requiring treatment.
  \item \textbf{Thyroid ultrasound with Doppler:} Heterogeneous hypoechoic texture, pseudonodules, and increased or decreased vascularity support Hashimoto thyroiditis. Used to evaluate for thyroid nodules requiring FNA (thyroid cancer risk is NOT increased by Hashimoto thyroiditis per se, though lymphoma risk is 60-fold higher).
  \item \textbf{Repeat TSH monitoring:} Annually for anti-TPO-positive euthyroid patients; every 6 months for subclinical hypothyroidism; every 4 weeks during pregnancy.
  \item \textbf{TRAb (second and third trimester):} In pregnant women with active or prior Graves disease, fetal thyroid status monitoring and neonatal risk assessment (TRAb >3x upper limit of normal indicates high risk).
  \item \textbf{Thyroglobulin with anti-Tg measurement:} In differentiated thyroid cancer follow-up. Anti-Tg status determines whether thyroglobulin is reliable as a biomarker.
  \item \textbf{Radioactive iodine uptake and scan:} In hyperthyroid patients with positive TRAb, confirms Graves disease etiology (diffuse increased uptake) and excludes subacute thyroiditis.
  \item \textbf{Fine-needle aspiration biopsy of thyroid nodules:} If thyroid nodules >1 cm are present in Hashimoto thyroiditis (TI-RADS guidelines apply, same as for non-autoimmune thyroid disease).
  \item \textbf{IgG4 subclass levels:} If IgG4-related thyroid disease is suspected (Riedel thyroiditis, fibrosing Hashimoto, IgG4-related hypophysitis).
\end{itemize}

\section*{References}
\begin{enumerate}
  \item Ross DS, Burch HB, Cooper DS, et al. 2016 American Thyroid Association guidelines for diagnosis and management of hyperthyroidism. Thyroid. 2016;26(10):1343-1421.
  \item Garber JR, Cobin RH, Gharib H, et al. Clinical practice guidelines for hypothyroidism in adults. Thyroid. 2012;22(12):1200-1235.
  \item Alexander EK, Pearce EN, Brent GA, et al. 2017 Guidelines of the American Thyroid Association for the diagnosis and management of thyroid disease during pregnancy and the postpartum. Thyroid. 2017;27(3):315-389.
  \item Vanderpump MP, Tunbridge WM, French JM, et al. The incidence of thyroid disorders in the community: a twenty-year follow-up of the Whickham Survey. Clin Endocrinol (Oxf). 1995;43(1):55-68.
  \item Spencer CA, Takeuchi M, Kazaros E, et al. Serum thyroglobulin autoantibodies: prevalence, influence on serum thyroglobulin measurement, and prognostic significance. J Clin Endocrinol Metab. 1998;83(4):1121-1127.
  \item Zophel K, Roggenbuck D, Schott M. Clinical review about TRAb assay's history. Autoimmun Rev. 2010;9(10):695-700.
  \item Baloch Z, Carayon P, Conte-Devolx B, et al. Laboratory medicine practice guidelines: laboratory support for the diagnosis and monitoring of thyroid disease. Thyroid. 2003;13(1):3-126.
  \item Burtis CA, Ashwood ER, Bruns DE. Tietz Textbook of Clinical Chemistry and Molecular Diagnostics. 6th ed. Elsevier; 2023.
\end{enumerate}

\clearpage
\section*{Regulatory Status and Authorization Date}
This chapter describes a test category, not a specific commercial product. FDA, CDSCO, EU IVDR, and MHRA approval or registration dates therefore cannot be assigned without the assay manufacturer, product name, instrument, and intended use. For laboratory-developed tests, an individual agency approval date may not exist. See the Preface for the interpretation of regulatory status.

\clearpage
